\chapter{Introduzione}
\label{cap:introduzione}

\section{Il contesto}

Il controllo di documenti aziendali è un processo laborioso e stancante. 
La verifica di ogni sezione del documento rispetto alla struttura e alle normative aziendali può essere un compito arduo, soprattutto per documenti di grandi dimensioni.
Per questo motivo, è nata l'idea di creare una piattaforma che faccia tutti questi controlli automaticamente tramite un modello AI esterno. 
Il modello, a sua volta, andrà ad effettuare controlli di qualità del documento e ritorerà una risposta strutturata in formato JSON al backend del sistema. La risposta è strutturata nel seguente modo:
\begin{itemize}
    \item Nome della sezione in cui si è verificato il problema
    \item I problemi rilevati e la soluzione proposta
    \item La gravità del problema
    \item La lista di sezioni che mancano
\end{itemize}
Grazie a queste informazioni, viene costruito un report dell'analisi che viene visualizzato nella dashboard e che può essere scaricato come file JSON oppure come PDF. Il progetto è stato proposto dall'azienda all'evento StageIT 2026.
\section{L'azienda}
Lo stage è stato svolto presso Var Group in Via Salboro 22, società operante nel settore dei servizi e delle soluzioni digitali fondata nel 1973 con sede legale ad Empoli.
L'azienda offre competenze in ambiti quali cloud computing, cybersecurity, data science, sviluppo software, intelligenza artificiale e trasformazione digitale, supportando organizzazioni pubbliche e private nell'innovazione dei propri processi tecnologici. \\


\section{Scelta dello stage}

Ho scelto questo progetto con l'obiettivo di approfondire le mie conoscenze nell'ambito dell'intelligenza artificiale.
La crescente diffusione dell'intelligenza artificiale ha rappresentato una delle principali motivazioni alla base della scelta di questo stage,
 offrendo l'opportunità di approfondire le potenzialità di tali tecnologie e comprenderne le applicazioni nello sviluppo di prodotti software.
\section{Organizzazione del documento}

Il documento è diviso in 7 capitoli, ognuno tratta una parte diversa del sistema e dell'esperienza dello stage.
La suddivisione avviene così:
\begin{itemize}
    \item Capitolo 1: introduzione del contenuto
    \item Capitolo 2: descrizione più profonda delle attività svolte durante lo stage
    \item Capitolo 3: descizione delle tecnologie utilizzate per realizzare il progetto
    \item Capitolo 4: descrizione delle scelte progettuali e architetturali
    \item Capitolo 5: descrizione della migrazione del sistema verso AWS cloud
    \item Capitolo 6: descrizione del processo di testing
    \item Capitolo 7: conclusione e valutazione personale dell'esperienza 
\end{itemize}


